\section{The fields, the roots, and the completed spaces}\label{rec:norms}

Our first objective is to define the norms to which the estimates apply.
All scalars are real. A rooted tree has a distinguished root, and every
node has finitely many predecessors. Its set of children is denoted by
$\operatorname{ch}(s)$. We allow countably many children.
For a finitely supported vector $v=(v_s)$, set the fields to zero below
its support and compute upwards:
\begin{align}
 A_s(v)&=\Bigl(\sum_{c\in\operatorname{ch}(s)}P_c(v)^2\Bigr)^{1/2},&
 B_s(v)&=\Bigl(\sum_{c\in\operatorname{ch}(s)}Q_c(v)^2\Bigr)^{1/2},\notag\\
 P_s(v)&=\max\{A_s(v),v_s+B_s(v)\},&
 Q_s(v)&=\max\{B_s(v),-v_s+A_s(v)\}.
 \label{rec:definition}
\end{align}
Only the finite predecessor closure of the support has nonzero fields.
Empty sums are zero. Thus this instruction is independent of the depth
at which the upward computation starts. At every node,
\begin{equation}\label{rec:field-properties}
 P_s(v)-Q_s(v)=v_s,\qquad P_s(v)\ge A_s(v),\qquad Q_s(v)\ge B_s(v),
\end{equation}
and at least one of the last two inequalities is an equality.
Indeed, subtracting $v_s$ from the formula for $P_s$ gives the formula
for $Q_s$. If the maximum for $P_s$ selects its first entry, its inequality
is tight; otherwise the inequality for $Q_s$ is tight.
Induction also shows that this is the coordinatewise least nonnegative
pair with the prescribed difference and child lower bounds.

\subsection{Three Euclidean constructions}

The local rule is the same in the following spaces, but the underlying
trees and the values taken at the root differ.
\begin{enumerate}[label=\textup{(\roman*)}]
\item On $T=\N^{<\omega}$, including the root $o=\varnothing$, put
\[
 N_\Sigma(v)=P_o(v)+Q_o(v),\qquad v\in c_{00}(T),
\]
and let $X_\Sigma$ be the completion.
\item Let $T_n=\bigcup_{j=0}^n\N^j$, $n\ge1$, be disjoint copies with
roots $r_n$. Add a root $\rho$ whose children are the $r_n$, obtaining a
tree $\mathcal J$. For vectors on $\mathcal J\setminus\{\rho\}$ set
$v_\rho=0$. Then $P_\rho(v)=Q_\rho(v)$; define
\[
 N_0(v)=P_\rho(v)=Q_\rho(v),
\]
and let $X_0$ be the completion.
\item\label{rec:joined-sum-definition}
On the same tree $\mathcal J$, now allowing the coordinate $v_\rho$ to
vary, put
\[
 N_J(v)=P_\rho(v)+Q_\rho(v),\qquad v\in c_{00}(\mathcal J),
\]
and let $X_J$ be the completion.
\end{enumerate}
In particular, $N_J(v)=2N_0(v)$ on the subspace where $v_\rho=0$.
Allowing that coordinate to vary produces an additional one-dimensional
part; it does not change the local recursion.

We call a finite root-containing set \emph{initial} if it contains all
predecessors of its nodes, write $\Pi_H$ for coordinate restriction,
and put $F_H=\ker\Pi_H$ in the corresponding completion.
The next theorem gives the principal quantitative conclusions. The definition
of $\widehat\delta$ is \eqref{intro:modulus}.

\begin{theorem}\label{rec:main}\label{rec:joined-root-stability}
For $N=N_\Sigma$ on $X_\Sigma$, or $N=N_J$ on $X_J$, one has
\begin{equation}\label{rec:sum-modulus}
 \widehat\delta_N(t)\ge t^3/128\qquad(0<t<1).
\end{equation}
For $X_0$, a finite initial set $H$, a vector $x$ supported in $H$, and
$z\in F_H\setminus\{0\}$, one has
\begin{equation}\label{rec:equal-estimate}
 \frac{N_0(x+z)+N_0(x-z)}2\ge N_0(x)
 +\frac{N_0(z)^3}{8(2N_0(x)+N_0(z))^2}.
\end{equation}
Each of $X_\Sigma,X_0,X_J$ admits no equivalent AUC norm.
\end{theorem}

The energy proof in Section~\ref{rec:energy} gives the first two
assertions, uniformly in the height and branching of $H$.
Section~\ref{rec:obstruction} proves the last assertion by following
bounded paths whose next increments are weakly null. The two later
probabilistic arguments obtain uniform separation by controlling the
propagation along random paths.

\begin{proposition}\label{rec:completion}
These formulas define infinite-dimensional Banach spaces after completion.
Coordinate evaluations and the fields extend continuously. If $H$ is a
finite set containing the root and all predecessors of its nodes, then
coordinate restriction $\Pi_H$ is a contractive projection. Its kernel
$F_H$ has finite codimension, and finitely supported elements of $F_H$
are dense in $F_H$. Both $X_0$ and $X_J$ are reflexive.
\end{proposition}

\begin{proof}[Proof of Proposition~\ref{rec:completion}]
The fields are positively homogeneous and subadditive. To verify
subadditivity, induct upwards: the child Euclidean norms are monotone
on nonnegative arrays and satisfy the triangle inequality. Each entry
in either maximum in \eqref{rec:definition} is therefore subadditive,
and a maximum of two subadditive functions is subadditive. Negation
interchanges $P_s$ and $Q_s$. This proves symmetry and the triangle
inequality for each proposed norm. Moreover,
\[
 |v_s|\le\max\{P_s(v),Q_s(v)\}\le N(v),
 \qquad N\in\{N_\Sigma,N_0,N_J\},
\]
by following the child inequalities to the corresponding root.
This proves definiteness and coordinate continuity. The coordinate
vectors are linearly independent, so the spaces are infinite dimensional.

The fields of $\Pi_Hv$ vanish outside $H$. Inside $H$ its coordinates
are unchanged, and each maximum is nondecreasing in the child fields;
induction gives $P_s(\Pi_Hv)\le P_s(v)$ and the same for $Q_s$.
Thus $\Pi_H$ is contractive. Ignore the fixed coordinate $\rho$ when
$N=N_0$. It extends to the completion and has finite-dimensional range.
Applying $I-\Pi_H$ to finite approximations proves tail density.
Choose an increasing sequence $(H_k)$ of finite initial sets covering
the tree. If a finite vector $v$ is supported in $H_k$, then
\[
 N(x-\Pi_{H_k}x)
 \le N(x-v)+N(\Pi_{H_k}(v-x))\le2N(x-v).
\]
Density therefore gives $\Pi_{H_k}x\to x$ for every completed vector
$x$. In particular, the coordinates separate points. A completed
vector whose coordinates vanish outside a finite initial set $H$
equals $\Pi_Hx$, so support in $H$ has its usual coordinate meaning.
Subadditivity and sign exchange give
\[
 |P_s(v)-P_s(w)|\le N(v-w),\qquad
 |Q_s(v)-Q_s(w)|\le N(v-w).
\]
The fields therefore extend continuously, and the root norm formulas
persist on the completions.

We prove reflexivity by describing the finite-height components. Set
$M_n(v)=\max\{P_{r_n}(v),Q_{r_n}(v)\}$ on $c_{00}(T_n)$, and let
$E_n$ be its completion under this norm.
Writing $m_s=\max\{P_s,Q_s\}$ gives
\[
 m_s\le |v_s|+\Bigl(\sum_{c\in\operatorname{ch}(s)}m_c^2\Bigr)^{1/2}.
\]
The Euclidean norm of the $m_s$ on one level is at most the Euclidean
norm of the coordinates on that level plus the corresponding norm of
$m_c$ on the next level. Iteration and Cauchy--Schwarz yield
$M_n(v)\le\sqrt{n+1}\,\|v\|_2$.
Conversely the $\ell_2$ norm of either field on any level is at most its
root value. Since $v=P-Q$,
\begin{equation}\label{rec:block-equivalence}
 \frac{\|v\|_2}{2\sqrt{n+1}}\le M_n(v)
 \le\sqrt{n+1}\,\|v\|_2.
\end{equation}
Hence $E_n$ is $\ell_2(T_n)$ with an equivalent norm,
and is reflexive. The recursion extends to its vectors by continuity
or directly by the finitely many levels.

For $v$ with zero joining coordinate, put
$L(v)=(\sum_n M_n(v|_{T_n})^2)^{1/2}$. The two child arrays at $\rho$
give
\begin{equation}\label{rec:joined-equivalence}
 N_0(v)\le L(v)\le\sqrt2\,N_0(v).
\end{equation}
For a free root coordinate $a=v_\rho$, with the remaining coordinates
written as $w$, let
$A=(\sum_nP_{r_n}(w)^2)^{1/2}$ and
$B=(\sum_nQ_{r_n}(w)^2)^{1/2}$.
Then
\begin{equation}\label{rec:free-root-equivalence}
 N_J(v)=\max\{2A-a,2B+a\},\qquad
 \max\{|a|,L(w)\}\le N_J(v)\le |a|+2L(w).
\end{equation}
For the lower bound, $N_J\ge |a|$ and $N_J\ge A+B\ge L(w)$.
These comparisons identify $X_0$ with the $\ell_2$ sum of the $E_n$
under an equivalent norm, and $X_J$ with that sum plus $\R$.
The Hilbert comparisons for individual components need not be uniform
in $n$: we take the outer sum in the norms $M_n$ themselves.

For completeness, an $\ell_2$ sum of Banach spaces is complete:
a Cauchy sequence converges in each component, and its Cauchy bounds
pass to all finite component sums and then to their supremum. Its dual
is the $\ell_2$ sum of the component duals. The upper norm bound is
Cauchy--Schwarz; the lower bound follows by testing a functional on
finitely many almost norming component vectors with scalar coefficients
proportional to their dual norms. Finite component vectors are dense.
Applying this dual identification twice proves reflexivity when each
component is reflexive. Equivalent norms preserve reflexivity.
This proves the assertion for both joined spaces.
\end{proof}

\paragraph{Component root sums.}
The alternative norm $P_{r_n}+Q_{r_n}$ lies between $M_n$ and $2M_n$,
so it gives the same component completion. Its behavior on individual
levels can also be read directly from the recursion. For a finite
vector on $T_n$, put $s_a=P_a+Q_a$. Then
\begin{equation}\label{rec:sum-level-bounds}
 \max\{|v_a|,\|(s_c)_{c\in\operatorname{ch}(a)}\|_2\}
 \le s_a\le |v_a|+2\|(s_c)_{c\in\operatorname{ch}(a)}\|_2.
\end{equation}
Indeed, $s_a\ge A_a+B_a\ge\|(s_c)_c\|_2$ and
$s_a=\max\{2A_a-v_a,2B_a+v_a\}$.
Iteration bounds its root value above by
$\sum_{d=0}^n2^d\|v|_{\text{level }d}\|_2$, while each level norm of
$v$ is at most that root value. These provide direct Hilbert comparisons
for the component root-sum convention.
